\paragraph{Mitigations}

Users of \putty are advised to update to version \puttyfixedversion or newer to avoid biased nonce generation when using ECDSA signatures with NIST~\pfivetwoone keys. Additionally, users should revoke any NIST~\pfivetwoone keys used with a \putty version below \puttyfixedversion by removing them from servers and platforms and regenerating them as necessary.

Implementers of cryptographic libraries need to carefully avoid biased nonces in (EC)DSA. As our attack on the deterministic nonce generation of \putty has shown, even deterministic nonce schemes may introduce bias. We recommend the algorithm described in~\cite{rfc6979} for the generation of deterministic nonces, which can generate nonces of arbitrary length.

Operators of online collaborative development platforms, and more generally server implementers, are advised to improve their public key security baselines in several aspects. First, they should test public keys on upload and reject keys based on the requirements outlined in~\cite{ca-browser-forum-2.1.2}. To complement this, platforms should perform periodic rescans of all public keys as necessary. While additional upload checks will necessarily add overhead, we estimate that most checks are computationally cheap and therefore reasonable to implement. Platform operators may decide to perform more expensive checks, such as batch GCD, offline.

Furthermore, platforms should reevaluate whether beneficial use cases\footnote{For example, the Ubuntu Linux installer can fetch the user's SSH public key from GitHub using only their GitHub username, providing a convenient way to install their public key on a new system, where copy\&paste may not be available.} for publishing their users' public keys outweigh the negative implications posed by privacy concerns, for example account linking, and scraping for weak keys.

On the SSH client side, vendors need to be aware of the design of the SSH agent interface. Obviously, it can be used by malicious servers to forge SSH signatures, potentially bypassing authentication guarantees. However, because it provides unrestricted access to a perfect signing oracle, it can also be vulnerable to attacks aiming for signature forgeries or private key extraction. Such a signing oracle may allow attackers to sign Git commits impersonating the victim, potentially bypassing repository restrictions. This may lead to unauthorized code or files being added to the repository and could facilitate supply chain attacks. When using modern continuous integration pipelines, this can also enable remote code execution on these systems when commits are built automatically. This can be somewhat mitigated by the agent checking the ``magic bytes'' that are present in Git commits and SSH client signatures. However, since this is a blocklist approach, it may only reduce the attack surface. Also, if side channels exist, for example, through timing, an attacker may use such an oracle to recover the private key used by the SSH agent. To mitigate both potential threats, \texttt{SSH\_AGENT\_CONSTRAIN\_CONFIRM} can be set by default for all keys loaded into an agent, requiring user interaction upon signature generation.

\paragraph{Impact of Weak Keys}

We discovered a variety of weak keys, including RSA keys with fewer than 830 bits and others that were easily factorable---some of which may appear to be merely dummy keys. However, regardless of the specific weakness, a weak key allows an attacker to authenticate itself against the SSH service of a Git-based platform as the affected user. This results in a direct relationship between the impact a weak key can potentially have and the repositories and corresponding privileges a user has access to. If, for example, a user uses the platform solely for private development projects, the impact may be straightforward. If, on the other hand, the user has write access to larger software projects or libraries, this can enable complex supply chain attacks by pushing inconspicuous commits to the repository. We did not explore the potential impact of the weak keys we found in greater detail, because finding all repositories a user has access to is non-trivial. For example, the user might have private repositories not visible to the public.
